\documentclass[a4paper,12pt]{paper}
%DIF LATEXDIFF DIFFERENCE FILE
%DIF DEL lecture9.tex             Thu Nov 13 02:11:28 2025
%DIF ADD diffs/lecture9_new.tex   Fri Sep 25 10:43:07 2026

\usepackage{mycommands}

\title{Lecture 20: Self-gravitation and rotation}
\author{David Al-Attar \\ Michaelmas Term, \the\year}
%DIF PREAMBLE EXTENSION ADDED BY LATEXDIFF
%DIF UNDERLINE PREAMBLE %DIF PREAMBLE
\RequirePackage[normalem]{ulem} %DIF PREAMBLE
\RequirePackage{color}\definecolor{RED}{rgb}{1,0,0}\definecolor{BLUE}{rgb}{0,0,1} %DIF PREAMBLE
\providecommand{\DIFadd}[1]{{\protect\color{blue}\uwave{#1}}} %DIF PREAMBLE
\providecommand{\DIFdel}[1]{{\protect\color{red}\sout{#1}}}                      %DIF PREAMBLE
%DIF SAFE PREAMBLE %DIF PREAMBLE
\providecommand{\DIFaddbegin}{} %DIF PREAMBLE
\providecommand{\DIFaddend}{} %DIF PREAMBLE
\providecommand{\DIFdelbegin}{} %DIF PREAMBLE
\providecommand{\DIFdelend}{} %DIF PREAMBLE
\providecommand{\DIFmodbegin}{} %DIF PREAMBLE
\providecommand{\DIFmodend}{} %DIF PREAMBLE
%DIF FLOATSAFE PREAMBLE %DIF PREAMBLE
\providecommand{\DIFaddFL}[1]{\DIFadd{#1}} %DIF PREAMBLE
\providecommand{\DIFdelFL}[1]{\DIFdel{#1}} %DIF PREAMBLE
\providecommand{\DIFaddbeginFL}{} %DIF PREAMBLE
\providecommand{\DIFaddendFL}{} %DIF PREAMBLE
\providecommand{\DIFdelbeginFL}{} %DIF PREAMBLE
\providecommand{\DIFdelendFL}{} %DIF PREAMBLE
%DIF COLORLISTINGS PREAMBLE %DIF PREAMBLE
\RequirePackage{listings} %DIF PREAMBLE
\RequirePackage{color} %DIF PREAMBLE
\lstdefinelanguage{DIFcode}{ %DIF PREAMBLE
%DIF DIFCODE_UNDERLINE %DIF PREAMBLE
  moredelim=[il][\color{red}\sout]{\%DIF\ <\ }, %DIF PREAMBLE
  moredelim=[il][\color{blue}\uwave]{\%DIF\ >\ } %DIF PREAMBLE
} %DIF PREAMBLE
\lstdefinestyle{DIFverbatimstyle}{ %DIF PREAMBLE
	language=DIFcode, %DIF PREAMBLE
	basicstyle=\ttfamily, %DIF PREAMBLE
	columns=fullflexible, %DIF PREAMBLE
	keepspaces=true %DIF PREAMBLE
} %DIF PREAMBLE
\lstnewenvironment{DIFverbatim}{\lstset{style=DIFverbatimstyle}}{} %DIF PREAMBLE
\lstnewenvironment{DIFverbatim*}{\lstset{style=DIFverbatimstyle,showspaces=true}}{} %DIF PREAMBLE
%DIF END PREAMBLE EXTENSION ADDED BY LATEXDIFF

\begin{document}

\maketitle
\DIFdelbegin %DIFDELCMD < 

%DIFDELCMD < %%%
\DIFdelend \section*{Outline and motivation}

In previous lectures we have discussed how \DIFdelbegin \DIFdel{that }\DIFdelend travel time and waveform observations can be used to determine variations in
 elastic wave speeds within the Earth. Nothing has, however, been said about \textbf{density}, with this parameter being key to 
 understanding both the Earth's interior composition and dynamics. To learn about density, we need to study the 
 \textbf{\DIFdelbegin \DIFdel{long period free-oscillations }\DIFdelend \DIFaddbegin \DIFadd{long-period free oscillations }\DIFaddend of the Earth}. This is because \textbf{self-gravitation}, and hence density, plays a vital role in 
 their physics. Within this lecture we will return to basic elastodynamics, and show how gravitational forces can be included 
 into the equations of motion. We also explain how the \textbf{Earth's rotation} enters into the problem. Having obtained the 
 equations of motion, we perform a \textbf{scaling analysis} to determine the relative importance of the different terms, and 
 conclude that at short enough periods the effects of self-gravitation and rotation can be safely neglected, and hence 
 the earlier parts of this course remain valid.

\section*{A review of Newtonian gravitation}

Consider the motion $\bphi$ of an elastic body relative to a reference body $M$, and write $M_{t}$ for 
the instantaneous volume of space occupied at time $t$. At this instant in time, the gravitational potential of the 
body satisfies Poisson\DIFdelbegin \DIFdel{equation
}\begin{displaymath}
\DIFdel{\nabla^{2}\phi=4\pi G\varrho.
}\end{displaymath}%DIFAUXCMD
\DIFdelend \DIFaddbegin \DIFadd{'s equation
}\begin{equation}
\DIFadd{\nabla^{2}\phi=4\pi G\varrho,
\label{eq:1}
}\end{equation}\DIFaddend 
in $\mathbb{R}^{3}$, where $G$ is the universal gravitational constant, and $\varrho$ is the density. The gravitational 
acceleration associated with this potential is given by
\DIFdelbegin \begin{displaymath}
\DIFdel{g_{i}=-\frac{\partial\phi}{\partial x_{i}}
}\end{displaymath}%DIFAUXCMD
\DIFdelend \DIFaddbegin \begin{equation}
\DIFadd{g_{i}=-\frac{\partial\phi}{\partial x_{i}}.
\label{eq:2}
}\end{equation}\DIFaddend 
\DIFdelbegin \DIFdel{The above Poisson}\DIFdelend \DIFaddbegin \DIFadd{Poisson's }\DIFaddend equation is subject to the conditions that $\phi$ tends to zero at infinity, while across any 
jumps in density we have the continuity conditions
\DIFdelbegin \begin{displaymath}
[\DIFdel{\phi}]\DIFdel{_{-}^{+}=0, \quad }[\DIFdel{g_{i}\hat{n}_{i}}]\DIFdel{_{-}^{+}=0
}\end{displaymath}%DIFAUXCMD
\DIFdelend \DIFaddbegin \begin{equation}
[\DIFadd{\phi}]\DIFadd{_{-}^{+}=0,\quad }[\DIFadd{g_{i}\hat{n}_{i}}]\DIFadd{_{-}^{+}=0,
\label{eq:3}
}\end{equation}\DIFaddend 
where $[\cdot]_{-}^{+}$ denotes the jump in a quantity in the direction of \DIFaddbegin \DIFadd{the }\DIFaddend outward normal.
The closed form solution to this problem is given by
\DIFdelbegin \begin{displaymath}
\DIFdel{\phi(\mathbf{y},t)=-G\int_{M_{t}}\frac{\varrho(\mathbf{y}',t)}{[(y_{i}-y_{i}')(y_{i}-y_{i}')]^{\frac{1}{2}}}\dd^{3}\mathbf{y}',
}\end{displaymath}%DIFAUXCMD
\DIFdelend \DIFaddbegin \begin{equation}
\DIFadd{\phi(\mathbf{y},t)=-G\int_{M_{t}}\frac{\varrho(\mathbf{y}',t)}{[(y_{i}-y_{i}')(y_{i}-y_{i}')]^{\frac{1}{2}}}\dd^{3}\mathbf{y}',
\label{eq:4}
}\end{equation}\DIFaddend 
and differentiating, we \DIFdelbegin \DIFdel{similar find
}\begin{displaymath}
\DIFdel{g_{i}(\mathbf{y},t)=-G\int_{M_{t}}\frac{\varrho(\mathbf{y}',t)(y_{i}-y_{i}')}{[(y_{j}-y_{j}')(y_{j}-y_{j}')]^{\frac{3}{2}}}\dd^{3}\mathbf{y}'.
}\end{displaymath}%DIFAUXCMD
\DIFdelend \DIFaddbegin \DIFadd{similarly find
}\begin{equation}
\DIFadd{g_{i}(\mathbf{y},t)=-G\int_{M_{t}}\frac{\varrho(\mathbf{y}',t)(y_{i}-y_{i}')}{[(y_{j}-y_{j}')(y_{j}-y_{j}')]^{\frac{3}{2}}}\dd^{3}\mathbf{y}'.
\label{eq:5}
}\end{equation}\DIFaddend 

It will be convenient to define the \textbf{referential gravitational potential} by
\DIFdelbegin \begin{displaymath}
\DIFdel{\zeta(\mathbf{x},t)=\phi[\bphi(\mathbf{x},t),t],
}\end{displaymath}%DIFAUXCMD
\DIFdelend \DIFaddbegin \begin{equation}
\DIFadd{\zeta(\mathbf{x},t)=\phi[\bphi(\mathbf{x},t),t],
\label{eq:6}
}\end{equation}\DIFaddend 
which is a time-dependent field on the reference body, $M$. We recall from Lecture\DIFdelbegin \DIFdel{12 }\DIFdelend \DIFaddbegin \DIFadd{~13 }\DIFaddend that the volume elements in 
the reference body and physical space are related by \DIFdelbegin \DIFdel{$\dd^{3}\mathbf{y} = J(\mathbf{x},t)\dd\mathbf{x}$ }\DIFdelend \DIFaddbegin \DIFadd{$\dd^{3}\mathbf{y} = J(\mathbf{x},t)\dd^{3}\mathbf{x}$ }\DIFaddend with $J=\det(\mathbf{F})$ 
the Jacobian of the motion. We can then change variables within eq.\DIFdelbegin \DIFdel{(4}\DIFdelend \DIFaddbegin \DIFadd{~(\ref{eq:4}}\DIFaddend ) and use the definition of $\zeta$ in eq.\DIFdelbegin \DIFdel{(6}\DIFdelend \DIFaddbegin \DIFadd{~(\ref{eq:6}}\DIFaddend ) to obtain
\DIFdelbegin \begin{displaymath}
\DIFdel{\zeta(\mathbf{x},t)=-G\int_{M}\frac{\rho(\mathbf{x}')}{\{[\varphi_{i}(\mathbf{x},t)-\varphi_{i}(\mathbf{x}',t)]
[\varphi_{i}(\mathbf{x},t)-\varphi_{i}(\mathbf{x}',t)]\}^{\frac{1}{2}}}\dd\mathbf{x}',
}\end{displaymath}%DIFAUXCMD
\DIFdelend \DIFaddbegin \begin{equation}
\DIFadd{\zeta(\mathbf{x},t)=-G\int_{M}\frac{\rho(\mathbf{x}')}{\{[\varphi_{i}(\mathbf{x},t)-\varphi_{i}(\mathbf{x}',t)]
[\varphi_{i}(\mathbf{x},t)-\varphi_{i}(\mathbf{x}',t)]\}^{\frac{1}{2}}}\dd^{3}\mathbf{x}',
\label{eq:7}
}\end{equation}\DIFaddend 
where from Lecture\DIFdelbegin \DIFdel{12 }\DIFdelend \DIFaddbegin \DIFadd{~13 }\DIFaddend we know that conservation of mass requires
\DIFdelbegin \begin{displaymath}
\DIFdel{\rho(\mathbf{x})=J(\mathbf{x},t)\varrho[\bphi(\mathbf{x},t),t],
}\end{displaymath}%DIFAUXCMD
\DIFdelend \DIFaddbegin \begin{equation}
\DIFadd{\rho(\mathbf{x})=J(\mathbf{x},t)\varrho[\bphi(\mathbf{x},t),t],
\label{eq:8}
}\end{equation}\DIFaddend 
with $\rho$ the \textbf{referential density}. In a similar manner, we define the \textbf{referential gravitational acceleration} by
\DIFdelbegin \begin{displaymath}
\DIFdel{\gamma_{i}(\mathbf{x},t)=g_{i}[\bphi(\mathbf{x},t),t]
}\end{displaymath}%DIFAUXCMD
\DIFdelend \DIFaddbegin \begin{equation}
\DIFadd{\gamma_{i}(\mathbf{x},t)=g_{i}[\bphi(\mathbf{x},t),t],
\label{eq:9}
}\end{equation}\DIFaddend 
which from eq.\DIFdelbegin \DIFdel{(5}\DIFdelend \DIFaddbegin \DIFadd{~(\ref{eq:5}}\DIFaddend ) can be written
\DIFdelbegin \begin{displaymath}
\DIFdel{\gamma_{i}(\mathbf{x},t)=-G\int_{M}\frac{\rho(\mathbf{x}')[\varphi_{i}(\mathbf{x},t)-\varphi_{i}(\mathbf{x}',t)]}{\{[\varphi_{j}(\mathbf{x},t)-\varphi_{j}(\mathbf{x}',t)][\varphi_{j}(\mathbf{x},t)-\varphi_{j}(\mathbf{x}',t)]\}^{\frac{3}{2}}}\dd\mathbf{x}',
}\end{displaymath}%DIFAUXCMD
\DIFdelend \DIFaddbegin \begin{equation}
\DIFadd{\gamma_{i}(\mathbf{x},t)=-G\int_{M}\frac{\rho(\mathbf{x}')[\varphi_{i}(\mathbf{x},t)-\varphi_{i}(\mathbf{x}',t)]}{\{[\varphi_{j}(\mathbf{x},t)-\varphi_{j}(\mathbf{x}',t)][\varphi_{j}(\mathbf{x},t)-\varphi_{j}(\mathbf{x}',t)]\}^{\frac{3}{2}}}\dd^{3}\mathbf{x}'.
\label{eq:10}
}\end{equation}\DIFaddend 

\subsection*{Gravitational binding energy}

At each instant of time we can associate a \textbf{gravitational binding energy} with the elastic body. This is the energy 
that would be required to assemble the body from matter dispersed at infinity, assuming that its constituent parts 
are acted on only by their mutual gravitational attraction. Because gravity is a conservative force, the binding 
energy is independent of the details of the assembly. 

Suppose that the mass of the body is already in place, 
and that its spatial density and gravitational potential are denoted by $\varrho$ and $\phi$. To determine an 
expression for the binding energy we will consider how this quantity changes due to an infinitesimal addition 
of mass to the body. By definition of the gravitational potential, the work required to bring in an additional 
mass element $\delta\varrho \dd^{3}\mathbf{y}$ from infinity to a point $\mathbf{y}$ within the body is
 $\delta\varrho\,\phi\dd^{3}\mathbf{y}$. 
The \DIFdelbegin \DIFdel{total }\DIFdelend \DIFaddbegin \DIFadd{change in the }\DIFaddend binding energy associated with a perturbation $\delta\varrho$ to \DIFaddbegin \DIFadd{the }\DIFaddend body's density \DIFdelbegin \DIFdel{by }\DIFdelend is, therefore, given by
\DIFdelbegin \begin{displaymath}
\DIFdel{\delta\mathcal{V}_{g}=\int_{M_{t}}\delta\varrho\,\phi \dd^{3}\mathbf{y}.
}\end{displaymath}%DIFAUXCMD
\DIFdelend \DIFaddbegin \begin{equation}
\DIFadd{\delta\mathcal{V}_{g}=\int_{M_{t}}\delta\varrho\,\phi \dd^{3}\mathbf{y},
\label{eq:11}
}\end{equation}\DIFaddend 
to first-order accuracy. Using the perturbed Poisson's equation $\nabla^{2}\delta\phi=4\pi G\delta\varrho$ for 
the resulting change in gravitational potential, we can write this increment to the binding energy as
\DIFdelbegin \begin{displaymath}
\DIFdel{\delta\mathcal{V}_{g}=\frac{1}{4\pi G}\int_{M_{t}}\left(\frac{\partial^{2}\delta\phi}{\partial x_{i}\partial x_{i}}\right)\phi
\dd^{3}\mathbf{y}=\frac{1}{4\pi G}\int_{\mathbb{R}^{3}}\left(\frac{\partial^{2}\delta\phi}{\partial x_{i}\partial x_{i}}\right)\phi\dd^{3}\mathbf{y},
}\end{displaymath}%DIFAUXCMD
\DIFdelend \DIFaddbegin \begin{equation}
\DIFadd{\delta\mathcal{V}_{g}=\frac{1}{4\pi G}\int_{M_{t}}\left(\frac{\partial^{2}\delta\phi}{\partial x_{i}\partial x_{i}}\right)\phi
\dd^{3}\mathbf{y}=\frac{1}{4\pi G}\int_{\mathbb{R}^{3}}\left(\frac{\partial^{2}\delta\phi}{\partial x_{i}\partial x_{i}}\right)\phi\dd^{3}\mathbf{y},
\label{eq:12}
}\end{equation}\DIFaddend 
where the final equality follows from the fact that $\nabla^{2}\delta\phi=0$ outside of the instantaneous 
body $M_{t}$. Using the identity
\DIFdelbegin \begin{displaymath}
\DIFdel{\left(\frac{\partial^{2}\delta\phi}{\partial x_{i}\partial x_{i}}\right)\phi=\frac{\partial}{\partial x_{i}}
\left(\phi\frac{\partial\delta\phi}{\partial x_{i}}\right)-\frac{\partial\phi}{\partial x_{i}}\frac{\partial\delta\phi}{\partial x_{i}}
}\end{displaymath}%DIFAUXCMD
\DIFdelend \DIFaddbegin \begin{equation}
\DIFadd{\left(\frac{\partial^{2}\delta\phi}{\partial x_{i}\partial x_{i}}\right)\phi=\frac{\partial}{\partial x_{i}}
\left(\phi\frac{\partial\delta\phi}{\partial x_{i}}\right)-\frac{\partial\phi}{\partial x_{i}}\frac{\partial\delta\phi}{\partial x_{i}},
\label{eq:13}
}\end{equation}\DIFaddend 
and applying the divergence theorem we obtain
\DIFdelbegin \begin{displaymath}
\DIFdel{\delta\mathcal{V}_{g}=-\frac{1}{4\pi G}\int_{\mathbb{R}^{3}}\frac{\partial\phi}{\partial x_{i}}\frac{\partial\delta\phi}{\partial x_{i}}\dd^{3}\mathbf{y}.
}\end{displaymath}%DIFAUXCMD
\DIFdelend \DIFaddbegin \begin{equation}
\DIFadd{\delta\mathcal{V}_{g}=-\frac{1}{4\pi G}\int_{\mathbb{R}^{3}}\frac{\partial\phi}{\partial x_{i}}\frac{\partial\delta\phi}{\partial x_{i}}\dd^{3}\mathbf{y},
\label{eq:14}
}\end{equation}\DIFaddend 
where we have used the continuity of the potentials and their normal derivatives across $\partial M_{t}$, while the 
surface integral at infinity vanishes because $\phi\sim r^{-1}$ and \DIFdelbegin \DIFdel{$||\nabla\delta\phi||\sim r^{-2}$ }\DIFdelend \DIFaddbegin \DIFadd{$\|\nabla\delta\phi\|\sim r^{-2}$ }\DIFaddend as $r\rightarrow\infty$. It follows that the increment to the gravitational binding energy can be written
\DIFdelbegin \begin{displaymath}
\DIFdel{\delta\mathcal{V}_{g}=\delta\left(-\frac{1}{8\pi G}\int_{\mathbb{R}^{3}}\frac{\partial\phi}{\partial x_{i}}
\frac{\partial\phi}{\partial x_{i}}\dd^{3}\mathbf{y}\right)
}\end{displaymath}%DIFAUXCMD
\DIFdelend \DIFaddbegin \begin{equation}
\DIFadd{\delta\mathcal{V}_{g}=\delta\left(-\frac{1}{8\pi G}\int_{\mathbb{R}^{3}}\frac{\partial\phi}{\partial x_{i}}
\frac{\partial\phi}{\partial x_{i}}\dd^{3}\mathbf{y}\right),
\label{eq:15}
}\end{equation}\DIFaddend 
and so we arrive at the result
\DIFdelbegin \begin{displaymath}
\DIFdel{\mathcal{V}_{g}=-\frac{1}{8\pi G}\int_{\mathbb{R}^{3}}\frac{\partial\phi}{\partial x_{i}}\frac{\partial\phi}{\partial x_{i}}\dd^{3}\mathbf{y}.
}\end{displaymath}%DIFAUXCMD
\DIFdelend \DIFaddbegin \begin{equation}
\DIFadd{\mathcal{V}_{g}=-\frac{1}{8\pi G}\int_{\mathbb{R}^{3}}\frac{\partial\phi}{\partial x_{i}}\frac{\partial\phi}{\partial x_{i}}\dd^{3}\mathbf{y}.
\label{eq:16}
}\end{equation}\DIFaddend 

Using \DIFdelbegin \DIFdel{the Poisson}\DIFdelend \DIFaddbegin \DIFadd{Poisson's }\DIFaddend equation for $\phi$\DIFaddbegin \DIFadd{, }\DIFaddend we can write the \DIFaddbegin \DIFadd{integrand of the }\DIFaddend above expression in the form
\DIFdelbegin \begin{displaymath}
\DIFdel{\frac{\partial\phi}{\partial x_{i}}\frac{\partial\phi}{\partial x_{i}}=\frac{\partial}{\partial x_{i}}
\left(\phi\frac{\partial\phi}{\partial x_{i}}\right)-\phi\left(\frac{\partial^{2}\phi}{\partial x_{i}
\partial x_{i}}\right)=\frac{\partial}{\partial x_{i}}\left(\phi\frac{\partial\phi}{\partial x_{i}}\right)-4\pi G\varrho\phi,
}\end{displaymath}%DIFAUXCMD
\DIFdelend \DIFaddbegin \begin{equation}
\DIFadd{\frac{\partial\phi}{\partial x_{i}}\frac{\partial\phi}{\partial x_{i}}=\frac{\partial}{\partial x_{i}}
\left(\phi\frac{\partial\phi}{\partial x_{i}}\right)-\phi\left(\frac{\partial^{2}\phi}{\partial x_{i}
\partial x_{i}}\right)=\frac{\partial}{\partial x_{i}}\left(\phi\frac{\partial\phi}{\partial x_{i}}\right)-4\pi G\varrho\phi,
\label{eq:17}
}\end{equation}\DIFaddend 
and applying the divergence theorem we obtain
\DIFdelbegin \begin{displaymath}
\DIFdel{\mathcal{V}_{g}(t)=\frac{1}{2}\int_{M_{t}}\varrho(\mathbf{y},t)\phi(\mathbf{y},t)\dd^{3}\mathbf{y}. ,
}\end{displaymath}%DIFAUXCMD
\DIFdelend \DIFaddbegin \begin{equation}
\DIFadd{\mathcal{V}_{g}(t)=\frac{1}{2}\int_{M_{t}}\varrho(\mathbf{y},t)\phi(\mathbf{y},t)\dd^{3}\mathbf{y},
\label{eq:18}
}\end{equation}\DIFaddend 
where, as above, we have used the continuity of $\phi$ and its normal derivative across $\partial M_{t}$, 
and the surface integral at infinity again vanishes. To include this binding energy into Hamilton's principle, 
it will be necessary to transform eq.\DIFdelbegin \DIFdel{(18}\DIFdelend \DIFaddbegin \DIFadd{~(\ref{eq:18}}\DIFaddend ) to the following integral over the reference body
\DIFdelbegin \begin{displaymath}
\DIFdel{\mathcal{V}_{g}(t)=\frac{1}{2}\int_{M}\rho(\mathbf{x})\zeta(\mathbf{x},t)\dd\mathbf{x}.
}\end{displaymath}%DIFAUXCMD
\DIFdelend \DIFaddbegin \begin{equation}
\DIFadd{\mathcal{V}_{g}(t)=\frac{1}{2}\int_{M}\rho(\mathbf{x})\zeta(\mathbf{x},t)\dd^{3}\mathbf{x},
\label{eq:19}
}\end{equation}\DIFaddend 
where we have used the definition of the referential gravitational potential in eq.\DIFdelbegin \DIFdel{(6}\DIFdelend \DIFaddbegin \DIFadd{~(\ref{eq:6}}\DIFaddend ) and the statement of 
conservation of mass during the deformation in eq.\DIFdelbegin \DIFdel{(8}\DIFdelend \DIFaddbegin \DIFadd{~(\ref{eq:8}}\DIFaddend ).

\subsection*{The equations of motion relative to inertial space}

Having obtained an appropriate expression for the gravitational binding energy, the equations of motion can 
be readily obtained using Hamilton's principle. The Lagrangian for the problem is
\DIFdelbegin \begin{displaymath}
\DIFdel{\mathcal{L}(\mathbf{x},t,\bphi,\mathbf{v},\mathbf{F})=\frac{1}{2}\rho(\mathbf{x})||\mathbf{v}||^{2}
-W(\mathbf{x},t,\mathbf{F})-\frac{1}{2}\rho(\mathbf{x})\zeta(\mathbf{x},t) ,
}\end{displaymath}%DIFAUXCMD
\DIFdelend \DIFaddbegin \begin{equation}
\DIFadd{\mathcal{L}(\mathbf{x},t,\bphi,\mathbf{v},\mathbf{F})=\frac{1}{2}\rho(\mathbf{x})\|\mathbf{v}\|^{2}
-W(\mathbf{x},t,\mathbf{F})-\frac{1}{2}\rho(\mathbf{x})\zeta(\mathbf{x},t),
\label{eq:20}
}\end{equation}\DIFaddend 
where we note that the elastic strain energy is allowed to be time-dependent to accommodate a seismic source. 
An important point about this Lagrangian is that it is \textbf{non-local} due to the gravitational binding energy. This 
means that the value of the Lagrangian at a point and time cannot be expressed only in terms of the values of 
the various fields and their derivatives at this point. The physical reason for this non-locality is, of course, 
the \textbf{\DIFdelbegin \DIFdel{action at a distance }\DIFdelend \DIFaddbegin \DIFadd{action-at-a-distance }\DIFaddend property of Newtonian gravitation}. 

The action for the body takes the usual form
\DIFdelbegin \begin{displaymath}
\DIFdel{\mathcal{S}(\bphi)=\int_{0}^{T}\int_{M}\mathcal{L}[\mathbf{x},t,\bphi
(\mathbf{x},t),\mathbf{v}(\mathbf{x},t),\mathbf{F}(\mathbf{x},t)]\dd\mathbf{x}\dd t,
}\end{displaymath}%DIFAUXCMD
\DIFdelend \DIFaddbegin \begin{equation}
\DIFadd{\mathcal{S}(\bphi)=\int_{0}^{T}\int_{M}\mathcal{L}[\mathbf{x},t,\bphi
(\mathbf{x},t),\mathbf{v}(\mathbf{x},t),\mathbf{F}(\mathbf{x},t)]\dd^{3}\mathbf{x}\dd t,
\label{eq:21}
}\end{equation}\DIFaddend 
and Hamilton's principle requires that this quantity be stationary about the true motion for all 
perturbations with fixed end points. In obtaining the \DIFdelbegin \DIFdel{Euler-Lagrange }\DIFdelend \DIFaddbegin \DIFadd{Euler--Lagrange }\DIFaddend equations, however, we have to account 
for the non-local nature of the Lagrangian. As ever, the mathematical statement of Hamilton's principle is 
\DIFdelbegin \DIFdel{that }\DIFdelend \DIFaddbegin \DIFadd{the }\DIFaddend vanishing of the functional derivative of the action
\DIFdelbegin \begin{displaymath}
\DIFdel{\langle D\mathcal{S}(\bphi)|\delta\bphi\rangle=0,
}\end{displaymath}%DIFAUXCMD
\DIFdelend \DIFaddbegin \begin{equation}
\DIFadd{\langle D\mathcal{S}(\bphi)|\delta\bphi\rangle=0,
\label{eq:22}
}\end{equation}\DIFaddend 
for all variations, $\delta\bphi$, subject to the usual fixed-endpoint conditions. Here the 
\DIFdelbegin \DIFdel{the left hand }\DIFdelend \DIFaddbegin \DIFadd{left-hand }\DIFaddend side is, by definition, just the linear term in the expansion of $\mathcal{S}(\bphi+
\delta\bphi)$ about $\bphi$. Looking at the form of the action, it is clear that this 
linear term takes the form
\DIFdelbegin \begin{displaymath}
\DIFdel{\langle D\mathcal{S}(\bphi)|\delta\bphi\rangle=\int_{0}^{T}\int_{M}
\left(\frac{\delta\mathcal{L}}{\delta\varphi_{i}}\delta\varphi_{i}+\frac{\delta\mathcal{L}}{\delta v_{i}}
\delta v_{i}+\frac{\delta\mathcal{L}}{\delta F_{ij}}\delta F_{ij}\right)\dd\mathbf{x}\dd t.
}\end{displaymath}%DIFAUXCMD
\DIFdelend \DIFaddbegin \begin{equation}
\DIFadd{\langle D\mathcal{S}(\bphi)|\delta\bphi\rangle=\int_{0}^{T}\int_{M}
\left(\frac{\delta\mathcal{L}}{\delta\varphi_{i}}\delta\varphi_{i}+\frac{\delta\mathcal{L}}{\delta v_{i}}
\delta v_{i}+\frac{\delta\mathcal{L}}{\delta F_{ij}}\delta F_{ij}\right)\dd^{3}\mathbf{x}\dd t,
\label{eq:23}
}\end{equation}\DIFaddend 
where we have introduced the convenient notation $\frac{\delta\mathcal{L}}{\delta\varphi_{i}}$, 
$\frac{\delta\mathcal{L}}{\delta v_{i}}$ and $\frac{\delta\mathcal{L}}{\delta F_{ij}}$ for the factors 
multiplying the various perturbations. In the case of a local Lagrangian these factors reduce to normal 
partial derivatives, and hence it is only the term $\frac{\delta\mathcal{L}}{\delta\varphi_{i}}$ that requires 
additional thought. Applying the usual integration by parts steps, we arrive at the \DIFdelbegin \DIFdel{Euler-Lagrange equation
}\begin{displaymath}
\DIFdel{\frac{\delta\mathcal{L}}{\delta\varphi_{i}}-\frac{\partial}{\partial t}\left(\frac{\delta\mathcal{L}}{\delta v_{i}}\right)-\frac{\partial}{\partial x_{j}}\left(\frac{\delta\mathcal{L}}{\delta F_{ij}}\right)=0,
}\end{displaymath}%DIFAUXCMD
\DIFdelend \DIFaddbegin \DIFadd{Euler--Lagrange equation
}\begin{equation}
\DIFadd{\frac{\delta\mathcal{L}}{\delta\varphi_{i}}-\frac{\partial}{\partial t}\left(\frac{\delta\mathcal{L}}{\delta v_{i}}\right)-\frac{\partial}{\partial x_{j}}\left(\frac{\delta\mathcal{L}}{\delta F_{ij}}\right)=0,
\label{eq:24}
}\end{equation}\DIFaddend 
along with the natural boundary condition
\DIFdelbegin \begin{displaymath}
\DIFdel{\frac{\delta\mathcal{L}}{\delta F_{ij}}\hat{n}_{j}=0
}\end{displaymath}%DIFAUXCMD
\DIFdelend \DIFaddbegin \begin{equation}
\DIFadd{\frac{\delta\mathcal{L}}{\delta F_{ij}}\hat{n}_{j}=0
\label{eq:25}
}\end{equation}\DIFaddend 
on $\partial M$. 

As in the non-gravitating case we have
\DIFdelbegin \begin{displaymath}
\DIFdel{\frac{\delta\mathcal{L}}{\delta v_{i}}=\rho v_{i}, \quad \frac{\delta\mathcal{L}}{\delta F_{ij}}=-\frac{\partial W}{\partial F_{ij}}
}\end{displaymath}%DIFAUXCMD
\DIFdelend \DIFaddbegin \begin{equation}
\DIFadd{\frac{\delta\mathcal{L}}{\delta v_{i}}=\rho v_{i},\quad \frac{\delta\mathcal{L}}{\delta F_{ij}}=-\frac{\partial W}{\partial F_{ij}},
\label{eq:26}
}\end{equation}\DIFaddend 
while in the \DIFdelbegin \DIFdel{final }\DIFdelend \DIFaddbegin \DIFadd{second }\DIFaddend problem set you will show that
\DIFdelbegin \begin{displaymath}
\DIFdel{\frac{\delta\mathcal{L}}{\delta\varphi_{i}}=\rho\gamma_{i}
}\end{displaymath}%DIFAUXCMD
\DIFdelend \DIFaddbegin \begin{equation}
\DIFadd{\frac{\delta\mathcal{L}}{\delta\varphi_{i}}=\rho\gamma_{i},
\label{eq:27}
}\end{equation}\DIFaddend 
with $\gamma_{i}$ the referential gravitational acceleration as defined above. The equations of motion, therefore, take the expected form
\DIFdelbegin \begin{displaymath}
\DIFdel{\rho\frac{\partial v_{i}}{\partial t}-\frac{\partial T_{ij}}{\partial x_{j}}-\rho\gamma_{i}=0,
}\end{displaymath}%DIFAUXCMD
\DIFdelend \DIFaddbegin \begin{equation}
\DIFadd{\rho\frac{\partial v_{i}}{\partial t}-\frac{\partial T_{ij}}{\partial x_{j}}-\rho\gamma_{i}=0,
\label{eq:28}
}\end{equation}\DIFaddend 
with $T_{ij}=\frac{\partial W}{\partial F_{ij}}$ the first \DIFdelbegin \DIFdel{Piola Kirchhoff }\DIFdelend \DIFaddbegin \DIFadd{Piola--Kirchhoff }\DIFaddend stress tensor. The form 
of these equations is not, of course, unexpected. To the balance of linear momentum obtained previously, 
we have simply added the gravitational force acting on the body due to its own mass.

\subsection*{Use of a co-rotating reference frame}
We have obtained the equations of motion for a self-gravitating body with respect to an \textbf{inertial reference} 
frame defined by the ``fixed stars'', and these results could be applied directly to the Earth. As, however, 
the Earth undergoes a nearly steady rotation about its polar axis, it is useful to reformulate the equations 
of motion with respect to a \textbf{non-inertial frame that rotates with the Earth}.

To do this we express the motion as
\DIFdelbegin \begin{displaymath}
\DIFdel{\varphi_{i}(\mathbf{x},t)=\Phi_{i}(t)+R_{ij}(t)\overline{\varphi}_{j}(\mathbf{x},t),
}\end{displaymath}%DIFAUXCMD
\DIFdelend \DIFaddbegin \begin{equation}
\DIFadd{\varphi_{i}(\mathbf{x},t)=\Phi_{i}(t)+R_{ij}(t)\overline{\varphi}_{j}(\mathbf{x},t),
\label{eq:29}
}\end{equation}\DIFaddend 
where $\mathbf{R}(t)$ is a rotation matrix. Here \DIFdelbegin \DIFdel{$\mathbf{\Phi}(t)$ denotes }\DIFdelend \DIFaddbegin \DIFadd{$\bm{\Phi}(t)$ denotes the }\DIFaddend translation of the reference frame relative 
to inertial space, and $\mathbf{R}(t)$ changes of orientation. The term $\overline{\bphi}(\mathbf{x},t)$ 
 describes the motion relative to the non-inertial frame, and will be called the \textbf{relative motion}.

 Differentiating eq.\DIFdelbegin \DIFdel{(29}\DIFdelend \DIFaddbegin \DIFadd{~(\ref{eq:29}}\DIFaddend ) with respect to time 
we see that the velocity relative to inertial space can be written
\DIFdelbegin \begin{displaymath}
\DIFdel{v_{i}=\frac{\ddns \Phi_{i}}{\ddns t}+R_{ij}\left(\frac{\partial\overline{\varphi}_{j}}{\partial t}+R_{kj}\frac{\ddns R_{kl}}{\ddns t}\overline{\varphi}_{l}\right)
}\end{displaymath}%DIFAUXCMD
\DIFdelend \DIFaddbegin \begin{equation}
\DIFadd{v_{i}=\frac{\ddns \Phi_{i}}{\ddns t}+R_{ij}\left(\frac{\partial\overline{\varphi}_{j}}{\partial t}+R_{kj}\frac{\ddns R_{kl}}{\ddns t}\overline{\varphi}_{l}\right),
\label{eq:30}
}\end{equation}\DIFaddend 
where we have used the identity $R_{kj}R_{kl}=\delta_{jl}$ for rotation matrices. 
In fact, if we differentiate this latter identity we find
\DIFdelbegin \begin{displaymath}
\DIFdel{\frac{\ddns R_{kj}}{\ddns t}R_{kl}+R_{kj}\frac{\ddns R_{kl}}{\ddns t}=0.
}\end{displaymath}%DIFAUXCMD
\DIFdelend \DIFaddbegin \begin{equation}
\DIFadd{\frac{\ddns R_{kj}}{\ddns t}R_{kl}+R_{kj}\frac{\ddns R_{kl}}{\ddns t}=0,
\label{eq:31}
}\end{equation}\DIFaddend 
and hence conclude that $R_{kj}\frac{\ddns R_{kl}}{\ddns t}$ is an \textbf{anti-symmetric matrix}. It follows that we can use the 
alternating tensor to write
\DIFdelbegin \begin{displaymath}
\DIFdel{R_{kj}\frac{\ddns R_{kl}}{\ddns t}=\epsilon_{jkl}\Omega_{k}
}\end{displaymath}%DIFAUXCMD
\DIFdelend \DIFaddbegin \begin{equation}
\DIFadd{R_{kj}\frac{\ddns R_{kl}}{\ddns t}=\epsilon_{jml}\Omega_{m},
\label{eq:32}
}\end{equation}\DIFaddend 
for a uniquely determined vector \DIFdelbegin \DIFdel{$\mathbf{\Omega}$}\DIFdelend \DIFaddbegin \DIFadd{$\bm{\Omega}$}\DIFaddend . Using this result, the velocity becomes
\DIFdelbegin \begin{displaymath}
\DIFdel{v_{i}=\frac{\ddns \Phi_{i}}{\ddns t}+R_{ij}(\overline{v}_{j}+\epsilon_{jkl}\Omega_{k}\overline{\varphi}_{l}) .
}\end{displaymath}%DIFAUXCMD
\DIFdelend \DIFaddbegin \begin{equation}
\DIFadd{v_{i}=\frac{\ddns \Phi_{i}}{\ddns t}+R_{ij}(\overline{v}_{j}+\epsilon_{jkl}\Omega_{k}\overline{\varphi}_{l}).
\label{eq:33}
}\end{equation}\DIFaddend 
Here $\overline{\mathbf{v}}$ is the velocity as seen from the non-inertial frame, while \DIFdelbegin \DIFdel{$\mathbf{\Omega}$ }\DIFdelend \DIFaddbegin \DIFadd{$\bm{\Omega}$ }\DIFaddend is the 
\textbf{angular velocity of the frame}. Differentiating again and applying the same ideas, the acceleration relative 
to inertial space can be written
\DIFdelbegin \begin{displaymath}
\DIFdel{\frac{\partial v_{i}}{\partial t}=\frac{\ddns ^{2}\Phi_{i}}{\ddns t^{2}}+R_{ij}\left(\frac{\partial\overline{v}_{j}}{\partial t}+2\epsilon_{jkl}\Omega_{k}\overline{v}_{l}+\epsilon_{jkl}\frac{\ddns \Omega_{k}}{\ddns t}\overline{\varphi}_{l}+\epsilon_{jkl}\epsilon_{lmn}\Omega_{k}\Omega_{m}\overline{\varphi}_{n}\right)
}\end{displaymath}%DIFAUXCMD
\DIFdelend \DIFaddbegin \begin{equation}
\DIFadd{\frac{\partial v_{i}}{\partial t}=\frac{\ddns ^{2}\Phi_{i}}{\ddns t^{2}}+R_{ij}\left(\frac{\partial\overline{v}_{j}}{\partial t}+2\epsilon_{jkl}\Omega_{k}\overline{v}_{l}+\epsilon_{jkl}\frac{\ddns \Omega_{k}}{\ddns t}\overline{\varphi}_{l}+\epsilon_{jkl}\epsilon_{lmn}\Omega_{k}\Omega_{m}\overline{\varphi}_{n}\right).
\label{eq:34}
}\end{equation}\DIFaddend 
On the \DIFdelbegin \DIFdel{right hand }\DIFdelend \DIFaddbegin \DIFadd{right-hand }\DIFaddend side we can identify various terms. The first is the acceleration of the frame relative to 
inertial space. Within the brackets we then see, relative to the non-inertial frame, accelerations associated 
with the motion along with the expected \textbf{Coriolis}, \textbf{Euler}, and \textbf{centrifugal} contributions.

We have not yet specified \emph{how} the non-inertial frame is to be chosen. Within seismology the most common 
approach is to take
\DIFdelbegin \begin{displaymath}
\DIFdel{\frac{\ddns \mathbf{\Phi}}{\ddns t}=\mathbf{0}, \quad \frac{\ddns \mathbf{\Omega}}{\ddns t}=\mathbf{0},
}\end{displaymath}%DIFAUXCMD
\DIFdelend \DIFaddbegin \begin{equation}
\DIFadd{\frac{\ddns \bm{\Phi}}{\ddns t}=\mathbf{0},\quad \frac{\ddns \bm{\Omega}}{\ddns t}=\mathbf{0},
\label{eq:35}
}\end{equation}\DIFaddend 
which corresponds to a frame whose \textbf{origin is static} relative to the \DIFdelbegin \DIFdel{"fixed stars" }\DIFdelend \DIFaddbegin \DIFadd{``fixed stars'' }\DIFaddend but \textbf{rotates steadily}, 
with the angular velocity chosen so that the frame \textbf{approximately co-rotates with the Earth}. The advantage of this 
choice is its simplicity. But it does mean that orbital, rotational, and relative motion of the Earth are mixed 
up within the term $\overline{\bphi}$. A preferable approach is to use the so-called 
\textbf{Tisserand frame} which is defined such that the relative motion, $\overline{\bphi}$ carries no net linear
 \DIFdelbegin \DIFdel{nor angular momentum}\footnote{\DIFdel{See ``On the elastodynamics of rotating planets'' by Maitra \& Al-Attar (2024) for details.}}%DIFAUXCMD
\addtocounter{footnote}{-1}%DIFAUXCMD
\DIFdel{.}\DIFdelend \DIFaddbegin \DIFadd{or angular momentum.}\footnote{\DIFadd{See Maitra \& Al-Attar (2024) for details.}}
\DIFaddend But we will stick to the simpler and more common choice in this course. 

To complete the derivation of the equations of motion in the non-inertial frame, we note from eq.\DIFdelbegin \DIFdel{(29) that
}\begin{displaymath}
\DIFdel{F_{ij}=R_{ik}\overline{F}_{kj}
}\end{displaymath}%DIFAUXCMD
\DIFdelend \DIFaddbegin \DIFadd{~(\ref{eq:29}) that
}\begin{equation}
\DIFadd{F_{ij}=R_{ik}\overline{F}_{kj},
\label{eq:36}
}\end{equation}\DIFaddend 
with $\overline{\mathbf{F}}$ the deformation gradient associated with relative motion. Due to material frame 
indifference, we then have $W(\mathbf{x},\mathbf{F})=W(\mathbf{x},\overline{\mathbf{F}})$ as should be physically 
expected. Using the chain rule, we find
\DIFdelbegin \begin{displaymath}
\DIFdel{\overline{T}_{ij}=\frac{\partial W}{\partial\overline{F}_{ij}}(\overline{\mathbf{F}})=\frac{\partial}{\partial
\overline{F}_{ij}}W(\mathbf{R}\overline{\mathbf{F}})=\frac{\partial W}{\partial F_{kl}}(\mathbf{R}\overline{\mathbf{F}})
\frac{\partial}{\partial\overline{F}_{ij}}(R_{km}\overline{F}_{ml})=R_{ki}T_{kj}
}\end{displaymath}%DIFAUXCMD
\DIFdelend \DIFaddbegin \begin{equation}
\DIFadd{\overline{T}_{ij}=\frac{\partial W}{\partial\overline{F}_{ij}}(\overline{\mathbf{F}})=\frac{\partial}{\partial
\overline{F}_{ij}}W(\mathbf{R}\overline{\mathbf{F}})=\frac{\partial W}{\partial F_{kl}}(\mathbf{R}\overline{\mathbf{F}})
\frac{\partial}{\partial\overline{F}_{ij}}(R_{km}\overline{F}_{ml})=R_{ki}T_{kj},
\label{eq:37}
}\end{equation}\DIFaddend 
with the \DIFdelbegin \DIFdel{left hand }\DIFdelend \DIFaddbegin \DIFadd{left-hand }\DIFaddend side the stress associated with $\overline{\bphi}$. Passing the rotation matrix to the 
other side, we have then obtained the useful identity
\DIFdelbegin \begin{displaymath}
\DIFdel{T_{ij}=R_{ik}\overline{T}_{kj}
}\end{displaymath}%DIFAUXCMD
\DIFdelend \DIFaddbegin \begin{equation}
\DIFadd{T_{ij}=R_{ik}\overline{T}_{kj},
\label{eq:38}
}\end{equation}\DIFaddend 
between the first \DIFdelbegin \DIFdel{Piola-Kirchhoff }\DIFdelend \DIFaddbegin \DIFadd{Piola--Kirchhoff }\DIFaddend stress relative to the inertial and non-inertial reference frames. Finally, we put 
eq.\DIFdelbegin \DIFdel{(29}\DIFdelend \DIFaddbegin \DIFadd{~(\ref{eq:29}}\DIFaddend ) into eq.\DIFdelbegin \DIFdel{(10}\DIFdelend \DIFaddbegin \DIFadd{~(\ref{eq:10}}\DIFaddend ) to find the relationship
\DIFdelbegin \begin{displaymath}
\DIFdel{\gamma_{i}=R_{ij}\overline{\gamma}_{i}
}\end{displaymath}%DIFAUXCMD
\DIFdelend \DIFaddbegin \begin{equation}
\DIFadd{\gamma_{i}=R_{ij}\overline{\gamma}_{j},
\label{eq:39}
}\end{equation}\DIFaddend 
between the referential gravitational acceleration in the two frames, with \DIFdelbegin \DIFdel{$\overline{\boldsymbol{\gamma}}$ }\DIFdelend \DIFaddbegin \DIFadd{$\overline{\bm{\gamma}}$ }\DIFaddend defined 
exactly as in eq.\DIFdelbegin \DIFdel{(10}\DIFdelend \DIFaddbegin \DIFadd{~(\ref{eq:10}}\DIFaddend ) but with $\overline{\bphi}$ replacing $\bphi$. Using these various 
results we have shown that the equations of motion are equivalent to
\DIFdelbegin \begin{displaymath}
\DIFdel{\rho\left(\frac{\partial\overline{v}_{i}}{\partial t}+2\epsilon_{ijk}\Omega_{j}\overline{v}_{k}+\epsilon_{ijk}
\epsilon_{klm}\Omega_{j}\Omega_{l}\overline{\varphi}_{m}\right)-\frac{\partial\overline{T}_{ij}}{\partial x_{j}}-\rho\overline{\gamma}_{i}=0.
}\end{displaymath}%DIFAUXCMD
\DIFdelend \DIFaddbegin \begin{equation}
\DIFadd{\rho\left(\frac{\partial\overline{v}_{i}}{\partial t}+2\epsilon_{ijk}\Omega_{j}\overline{v}_{k}+\epsilon_{ijk}
\epsilon_{klm}\Omega_{j}\Omega_{l}\overline{\varphi}_{m}\right)-\frac{\partial\overline{T}_{ij}}{\partial x_{j}}-\rho\overline{\gamma}_{i}=0.
\label{eq:40}
}\end{equation}\DIFaddend 
These equations are identical in form to those in an inertial frame except for the inclusion of the Coriolis and 
centrifugal accelerations -- the Euler acceleration has been removed because we have restricted attention to a 
steadily rotating frame. From this point on we will consistently work in the approximately co-rotating frame, 
and hence will simplify notations by dropping overlines on terms linked to the relative motion. 

\subsection*{Linearised equations of motion}

As ever, our main concern is with the linearised equations of motion about an equilibrium state.
Strictly, we should use the term \textbf{relative equilibrium}, as the body is now 
undergoing a \DIFdelbegin \DIFdel{steadily }\DIFdelend \DIFaddbegin \DIFadd{steady }\DIFaddend rotation relative to inertial space. We consider a disturbance to the body
caused by a stress glut, and hence \DIFaddbegin \DIFadd{take }\DIFaddend the strain energy to be
\DIFdelbegin \begin{displaymath}
\DIFdel{W(\mathbf{x},t,\mathbf{F})=U(\mathbf{x},\mathbf{C})-s\frac{1}{2}\overline{S}_{ij}(\mathbf{x},t)C_{ij}
}\end{displaymath}%DIFAUXCMD
\DIFdelend \DIFaddbegin \begin{equation}
\DIFadd{W(\mathbf{x},t,\mathbf{F})=U(\mathbf{x},\mathbf{C})-s\frac{1}{2}\overline{S}_{ij}(\mathbf{x},t)C_{ij},
\label{eq:41}
}\end{equation}\DIFaddend 
with $s$ a perturbation parameter, and $\overline{\mathbf{S}}$ a given stress glut. \DIFaddbegin \DIFadd{As in Lecture~15, the elastic tensor 
that appears below is the second derivative of the time-independent part $U(\mathbf{x},\mathbf{C})$ of this strain energy function. 
}\DIFaddend In a now standard manner we can 
obtain linearised equations of motion for the displacement vector defined as the first-order term in the expansion
\DIFdelbegin \begin{displaymath}
\DIFdel{\varphi_{i}(\mathbf{x},t)=x_{i}+s\,u_{i}(\mathbf{x},t)+O(s^{2}),
}\end{displaymath}%DIFAUXCMD
\DIFdelend \DIFaddbegin \begin{equation}
\DIFadd{\varphi_{i}(\mathbf{x},t)=x_{i}+s\,u_{i}(\mathbf{x},t)+O(s^{2}).
\label{eq:42}
}\end{equation}\DIFaddend 
At \DIFdelbegin \DIFdel{zeroth-order }\DIFdelend \DIFaddbegin \DIFadd{zeroth order, }\DIFaddend the equations of motion take the form
\DIFdelbegin \begin{displaymath}
\DIFdel{\rho\epsilon_{ijk}\epsilon_{klm}\Omega_{j}\Omega_{l}x_{m}-\frac{\partial T_{ij}^{0}}{\partial x_{j}}-\rho\gamma_{i}^{0}=0,
}\end{displaymath}%DIFAUXCMD
\DIFdelend \DIFaddbegin \begin{equation}
\DIFadd{\rho\epsilon_{ijk}\epsilon_{klm}\Omega_{j}\Omega_{l}x_{m}-\frac{\partial T_{ij}^{0}}{\partial x_{j}}-\rho\gamma_{i}^{0}=0,
\label{eq:43}
}\end{equation}\DIFaddend 
where $T_{ij}^{0}$ and $\gamma_{i}^{0}$ denote the equilibrium values of the stress tensor and gravitational acceleration. 
We will have more to say about these equations in the following lecture, but for the moment note that there 
must be a balance between the centrifugal \DIFdelbegin \DIFdel{, }\DIFdelend \DIFaddbegin \DIFadd{and }\DIFaddend gravitational forces on the body and \DIFdelbegin \DIFdel{a }\DIFdelend \DIFaddbegin \DIFadd{the }\DIFaddend resulting equilibrium stress; the 
simplifying assumption of a stress-free equilibrium is no longer an option. 

From the first-order term in the expansion we arrive at the linearised equations of motion for the displacement field
\DIFdelbegin \begin{displaymath}
\DIFdel{\rho\left(\frac{\partial^{2}u_{i}}{\partial t^{2}}+2\epsilon_{ijk}\Omega_{j}\frac{\partial u_{k}}{\partial t}+
\epsilon_{ijk}\epsilon_{klm}\Omega_{j}\Omega_{l}u_{m}\right)-\frac{\partial}{\partial x_{j}}\left(A_{ijkl}
\frac{\partial u_{k}}{\partial x_{l}}\right)-\rho\gamma_{i}^{1}=-\frac{\partial\overline{S}_{ij}}{\partial x_{j}},
}\end{displaymath}%DIFAUXCMD
\DIFdelend \DIFaddbegin \begin{equation}
\DIFadd{\rho\left(\frac{\partial^{2}u_{i}}{\partial t^{2}}+2\epsilon_{ijk}\Omega_{j}\frac{\partial u_{k}}{\partial t}+
\epsilon_{ijk}\epsilon_{klm}\Omega_{j}\Omega_{l}u_{m}\right)-\frac{\partial}{\partial x_{j}}\left(A_{ijkl}
\frac{\partial u_{k}}{\partial x_{l}}\right)-\rho\gamma_{i}^{1}=-\frac{\partial\overline{S}_{ij}}{\partial x_{j}},
\label{eq:44}
}\end{equation}\DIFaddend 
where $\gamma_{i}^{1}$ denotes the first-order perturbation to the referential gravitational \DIFdelbegin \DIFdel{potential}\DIFdelend \DIFaddbegin \DIFadd{acceleration}\DIFaddend . 
Using a binomial expansion, this term can be written explicitly as
\DIFdelbegin \begin{displaymath}
\DIFdel{\gamma_{i}^{1}=-G\int_{M}\rho'\left(\frac{\delta_{ij}}{[(x_{k}-x_{k}')(x_{k}-x_{k}')]^{3/2}}-
\frac{3(x_{i}-x_{i}')(x_{j}-x_{j}')}{[(x_{k}-x_{k}')(x_{k}-x_{k}')]^{5/2}}\right)(u_{j}-u_{j}')\dd\mathbf{x}',
}\end{displaymath}%DIFAUXCMD
\DIFdelend \DIFaddbegin \begin{equation}
\DIFadd{\gamma_{i}^{1}=-G\int_{M}\rho'\left(\frac{\delta_{ij}}{[(x_{k}-x_{k}')(x_{k}-x_{k}')]^{3/2}}-
\frac{3(x_{i}-x_{i}')(x_{j}-x_{j}')}{[(x_{k}-x_{k}')(x_{k}-x_{k}')]^{5/2}}\right)(u_{j}-u_{j}')\dd^{3}\mathbf{x}',
\label{eq:45}
}\end{equation}\DIFaddend 
where it is understood that primed terms are functions of the integration variable. The boundary conditions 
for the problem take the expected form
\DIFdelbegin \begin{displaymath}
\DIFdel{A_{ijkl}\frac{\partial u_{k}}{\partial x_{l}}\hat{n}_{j}=\overline{S}_{ij}\hat{n}_{j} .
}\end{displaymath}%DIFAUXCMD
\DIFdelend \DIFaddbegin \begin{equation}
\DIFadd{A_{ijkl}\frac{\partial u_{k}}{\partial x_{l}}\hat{n}_{j}=\overline{S}_{ij}\hat{n}_{j}.
\label{eq:46}
}\end{equation}\DIFaddend 
The key point is that the inclusion of gravitation and rotation into the problem only requires some additional 
terms in the equations of motion considered so far. It is notable, however, that these are no longer partial 
differential equations, but \textbf{\DIFdelbegin \DIFdel{integro partial differential }\DIFdelend \DIFaddbegin \DIFadd{integro-differential }\DIFaddend equations}, with 
both derivatives and integrals of the unknown functions occurring within the equations of motion. 
Again, this feature derives from the \DIFdelbegin \DIFdel{action at a distance }\DIFdelend \DIFaddbegin \DIFadd{action-at-a-distance }\DIFaddend property of Newtonian gravitation.

\subsection*{The neglect of self-gravitation and rotation}

We have obtained the linearised equations of motion within a steadily rotating and self-gravitating body. 
It is natural to ask if and why it has been reasonable to neglect these physical phenomena up till this point. 
To do so, we will perform a scaling analysis of the equations. 

We assume that the displacement vector field is 
characterised by a typical amplitude \DIFdelbegin \DIFdel{$||\mathbf{u}||\sim U$, a length-scale }\DIFdelend \DIFaddbegin \DIFadd{$\|\mathbf{u}\|\sim U$, a length scale }\DIFaddend $L$ and \DIFdelbegin \DIFdel{time-scale }\DIFdelend \DIFaddbegin \DIFadd{a time scale }\DIFaddend $T$. The following 
\DIFdelbegin \DIFdel{order of magnitude }\DIFdelend \DIFaddbegin \DIFadd{order-of-magnitude }\DIFaddend scaling estimates are then immediate
\DIFdelbegin \begin{eqnarray*}
\DIFdel{\rho\frac{\partial^{2}u_{i}}{\partial t^{2}}}&\DIFdel{\sim\frac{\overline{\rho}U}{T^{2}}, }\\
\DIFdel{\rho\epsilon_{ijk}\Omega_{j}\frac{\partial u_{k}}{\partial t}}&\DIFdel{\sim\frac{\overline{\rho}\Omega U}{T}; }\\
\DIFdel{\rho\epsilon_{ijk}\epsilon_{klm}\Omega_{j}\Omega_{l}u_{m}}&\DIFdel{\sim2\overline{\rho}\Omega^{2}U, }\\
\DIFdel{\frac{\partial}{\partial x_{j}}\left(A_{ijkl}\frac{\partial u_{k}}{\partial x_{l}}\right)}&\DIFdel{\sim\frac{\overline{\mu}U}{L^{2}},
}\end{eqnarray*}%DIFAUXCMD
\DIFdelend \DIFaddbegin \begin{align}
\DIFadd{\rho\frac{\partial^{2}u_{i}}{\partial t^{2}}}&\DIFadd{\sim\frac{\overline{\rho}U}{T^{2}}, \label{eq:47} }\\
\DIFadd{\rho\epsilon_{ijk}\Omega_{j}\frac{\partial u_{k}}{\partial t}}&\DIFadd{\sim\frac{\overline{\rho}\Omega U}{T}, \label{eq:48} }\\
\DIFadd{\rho\epsilon_{ijk}\epsilon_{klm}\Omega_{j}\Omega_{l}u_{m}}&\DIFadd{\sim2\overline{\rho}\Omega^{2}U, \label{eq:49} }\\
\DIFadd{\frac{\partial}{\partial x_{j}}\left(A_{ijkl}\frac{\partial u_{k}}{\partial x_{l}}\right)}&\DIFadd{\sim\frac{\overline{\mu}U}{L^{2}}, \label{eq:50}
}\end{align}\DIFaddend 
where $\overline{\rho}$ and $\overline{\mu}$ denote typical values of density and elastic \DIFdelbegin \DIFdel{modulii }\DIFdelend \DIFaddbegin \DIFadd{moduli }\DIFaddend within the body, 
while $\Omega$ is the magnitude of the rotation vector. The scaling estimate for the gravitational term within 
the equations of motion requires more thought. The referential gravitational potential at equilibrium is given by
\DIFdelbegin \begin{displaymath}
\DIFdel{\zeta_{0}=-G\int_{M}\frac{\rho'}{[(x_{k}-x_{k}')(x_{k}-x_{k}')]^{1/2}}\dd\mathbf{x}'.
}\end{displaymath}%DIFAUXCMD
\DIFdelend \DIFaddbegin \begin{equation}
\DIFadd{\zeta_{0}=-G\int_{M}\frac{\rho'}{[(x_{k}-x_{k}')(x_{k}-x_{k}')]^{1/2}}\dd^{3}\mathbf{x}'.
\label{eq:51}
}\end{equation}\DIFaddend 
Taking the gradient of this expression twice, we obtain
\DIFdelbegin \begin{displaymath}
\DIFdel{\frac{\partial\zeta_{0}}{\partial x_{i}\partial x_{j}}=G\int_{M}\rho'\left(\frac{\delta_{ij}}{[(x_{k}-x_{k}')
(x_{k}-x_{k}')]^{3/2}}-\frac{3(x_{i}-x_{i}')(x_{j}-x_{j}')}{[(x_{k}-x_{k}')(x_{k}-x_{k}')]^{5/2}}\right)\dd\mathbf{x}',
}\end{displaymath}%DIFAUXCMD
\DIFdelend \DIFaddbegin \begin{equation}
\DIFadd{\frac{\partial^{2}\zeta_{0}}{\partial x_{i}\partial x_{j}}=G\int_{M}\rho'\left(\frac{\delta_{ij}}{[(x_{k}-x_{k}')
(x_{k}-x_{k}')]^{3/2}}-\frac{3(x_{i}-x_{i}')(x_{j}-x_{j}')}{[(x_{k}-x_{k}')(x_{k}-x_{k}')]^{5/2}}\right)\dd^{3}\mathbf{x}',
\label{eq:52}
}\end{equation}\DIFaddend 
and using \DIFdelbegin \DIFdel{the Poisson}\DIFdelend \DIFaddbegin \DIFadd{Poisson's }\DIFaddend equation $\nabla^{2}\zeta_{0}=4\pi G\rho$ along with eq.\DIFdelbegin \DIFdel{(45}\DIFdelend \DIFaddbegin \DIFadd{~(\ref{eq:45}}\DIFaddend ) we arrive at the desired scaling estimate
\DIFdelbegin \begin{displaymath}
\DIFdel{\rho\gamma_{i}^{1}\sim4\pi G\overline{\rho}^{2}U
}\end{displaymath}%DIFAUXCMD
\DIFdelend \DIFaddbegin \begin{equation}
\DIFadd{\rho\gamma_{i}^{1}\sim4\pi G\overline{\rho}^{2}U.
\label{eq:53}
}\end{equation}\DIFaddend 

Using these results, we can quantify the relative importance of the terms within the force balance. First we have
\DIFdelbegin \begin{displaymath}
\DIFdel{\frac{\text{gravitational forces}}{\text{elastic forces}}=\frac{4\pi G\overline{\rho}^{2}U}{\frac{\overline{\mu} U}{L^{2}}}=\frac{4\pi G\overline{\rho}L^{2}}{\overline{v}^{2}},
}\end{displaymath}%DIFAUXCMD
\DIFdelend \DIFaddbegin \begin{equation}
\DIFadd{\frac{\text{gravitational forces}}{\text{elastic forces}}=\frac{4\pi G\overline{\rho}^{2}U}{\frac{\overline{\mu} U}{L^{2}}}=\frac{4\pi G\overline{\rho}L^{2}}{\overline{v}^{2}},
\label{eq:54}
}\end{equation}\DIFaddend 
where \DIFdelbegin \DIFdel{$\overline{v}=\sqrt{\frac{\overline{\mu}}{\rho}}$ }\DIFdelend \DIFaddbegin \DIFadd{$\overline{v}=\sqrt{\overline{\mu}/\overline{\rho}}$ }\DIFaddend represents a typical elastic wave speed in the model. 
This ratio scales like $L^{2}$ where $L$ is the \DIFdelbegin \DIFdel{characteristics length-scale }\DIFdelend \DIFaddbegin \DIFadd{characteristic length scale }\DIFaddend of the deformation. We conclude that 
for deformations occurring over sufficiently small \DIFdelbegin \DIFdel{length-scales}\DIFdelend \DIFaddbegin \DIFadd{length scales}\DIFaddend , gravitational forces will be negligible relative 
to elastic ones. We should then ask whether gravitational forces will ever be of importance within the Earth\DIFdelbegin \DIFdel{? }\DIFdelend \DIFaddbegin \DIFadd{. }\DIFaddend To 
answer this, we note that the largest-scale deformation within the Earth can have $L$ equal to its average radius \DIFdelbegin \DIFdel{6371 km}\DIFdelend \DIFaddbegin \DIFadd{$6371\,\mathrm{km}$}\DIFaddend .
Recalling that \DIFdelbegin \DIFdel{$G=6.67\times10^{-11} \, \text{m}^{3}\text{kg}^{-1}\text{s}^{-2}$}\DIFdelend \DIFaddbegin \DIFadd{$G=6.67\times10^{-11}\,\mathrm{m^{3}\,kg^{-1}\,s^{-2}}$}\DIFaddend , and taking as typical values \DIFdelbegin \DIFdel{$\overline{\rho}=5000 \, 
\text{kg m}^{-3}$ and $\overline{v}=8000 \, \text{ms}^{-1}$ we find
}\begin{displaymath}
\DIFdel{\frac{\text{gravitational forces}}{\text{elastic forces}}\sim3.
}\end{displaymath}%DIFAUXCMD
\DIFdelend \DIFaddbegin \DIFadd{$\overline{\rho}=5000\,
\mathrm{kg\,m^{-3}}$ and $\overline{v}=8000\,\mathrm{m\,s^{-1}}$, we find
}\begin{equation}
\DIFadd{\frac{\text{gravitational forces}}{\text{elastic forces}}\sim3.
\label{eq:55}
}\end{equation}\DIFaddend 
For such planetary scale deformations, it follows that gravitational forces are of the same order of magnitude as elastic forces, 
and so cannot be neglected within the dynamics. 

Next, we consider the ratio of the Coriolis and inertial terms, and find
\DIFdelbegin \begin{displaymath}
\DIFdel{\frac{\text{Coriolis term}}{\text{inertial term}}=\frac{\frac{\overline{\rho}\Omega U}{T}}{\frac{\overline{\rho}U}{T^{2}}}=\Omega T
}\end{displaymath}%DIFAUXCMD
\DIFdelend \DIFaddbegin \begin{equation}
\DIFadd{\frac{\text{Coriolis term}}{\text{inertial term}}=\frac{\frac{\overline{\rho}\Omega U}{T}}{\frac{\overline{\rho}U}{T^{2}}}=\Omega T.
\label{eq:56}
}\end{equation}\DIFaddend 
It follows that the Coriolis terms will be negligible for motions whose \DIFdelbegin \DIFdel{time-scale }\DIFdelend \DIFaddbegin \DIFadd{time scale }\DIFaddend is short relative to the rotation period 
of the Earth. A similar conclusion is readily shown to hold for the centrifugal term. It is known that seismic deformation 
is limited to periods shorter than about one hour, and so rotational effects play a relatively small (but not negligible) 
role \DIFaddbegin \DIFadd{in }\DIFaddend the problem. There are, however, processes at longer periods, e.g.\DIFaddbegin \DIFadd{\ }\DIFaddend those associated with solid-Earth tides, for which 
rotation becomes very important.

\section*{What you need to know and be able to do}
\DIFdelbegin %DIFDELCMD < \begin{itemize}
%DIFDELCMD <     %%%
\DIFdelend \DIFaddbegin \begin{enumerate}
    \DIFaddend \item[(i)] Understand the relevant parts of Newtonian gravity, including the manner in which the gravitational potential 
    and acceleration can be transformed into quantities defined on the reference body.
    \item[(ii)] Know physically what the gravitational binding energy of a planet represents, and how it can be determined.
    \item[(iii)] Know how the gravitational binding energy is incorporated into the Lagrangian, and how the equations of 
    motion can be obtained. The full derivation is too long to be examined, but you may be asked to reproduce parts of it 
    given appropriate information and guidance.
    \item[(iv)] Perform scaling arguments to determine the relative size of terms \DIFdelbegin \DIFdel{with in }\DIFdelend \DIFaddbegin \DIFadd{within }\DIFaddend an equation, and to determine 
    whether some of them might be negligible. This is a technique required throughout this course, and you should \DIFdelbegin \DIFdel{insure }\DIFdelend \DIFaddbegin \DIFadd{ensure }\DIFaddend you are happy with it.
\DIFdelbegin %DIFDELCMD < \end{itemize}
%DIFDELCMD < %%%
\DIFdelend \DIFaddbegin \end{enumerate}
\DIFaddend 

\end{document}
